\subsubsection{Conceptos y formulas de geometria 2D}

% Guia compacta de geometria analitica y vectorial para CP

\paragraph{1. Vectores, Dot Product y Cross Product.}
Un punto o vector en 2D es $P = (x, y)$. La resta $\vec{u} = B - A = (x_B - x_A, y_B - y_A)$ representa el vector que va de $A$ a $B$.
\begin{itemize}\setlength\itemsep{0pt}
	\item \textbf{Magnitud}: $|\vec{u}| = \sqrt{x^2 + y^2}$, con $|\vec{u}|^2 = x^2 + y^2$. En CP, \textbf{siempre que sea posible compara distancias al cuadrado} para trabajar con enteros y evitar perdida de precision por punto flotante.
	\item \textbf{Producto Punto (Dot Product)}:
	\[
		\vec{u} \cdot \vec{v} = x_1 x_2 + y_1 y_2 = |\vec{u}||\vec{v}| \cos \theta
	\]
	\begin{itemize}\setlength\itemsep{0pt}
		\item $\vec{u} \cdot \vec{v} > 0 \iff \theta < 90^\circ$ (angulo agudo).
		\item $\vec{u} \cdot \vec{v} = 0 \iff \theta = 90^\circ$ (\textbf{perpendiculares / ortogonales}).
		\item $\vec{u} \cdot \vec{v} < 0 \iff \theta > 90^\circ$ (angulo obtuso).
		\item Proyeccion escalar de $\vec{u}$ sobre $\vec{v}$: $\frac{\vec{u} \cdot \vec{v}}{|\vec{v}|}$. Proyeccion vectorial: $\text{proj}_{\vec{v}}(\vec{u}) = \frac{\vec{u} \cdot \vec{v}}{|\vec{v}|^2} \vec{v}$.
	\end{itemize}
	\item \textbf{Producto Cruz 2D (Cross Product / Pseudocruz)}:
	\[
		\vec{u} \times \vec{v} = x_1 y_2 - y_1 x_2 = |\vec{u}||\vec{v}| \sin \theta
	\]
	\begin{itemize}\setlength\itemsep{0pt}
		\item $|\vec{u} \times \vec{v}|$ = \textbf{Area del paralelogramo} formado por $\vec{u}$ y $\vec{v}$.
		\item Area con signo del triangulo $\triangle ABC$: $\frac{1}{2}((B - A) \times (C - A))$.
		\item $\vec{u} \times \vec{v} > 0 \iff \vec{v}$ esta a la \textbf{izquierda} de $\vec{u}$ (giro antihorario / CCW).
		\item $\vec{u} \times \vec{v} < 0 \iff \vec{v}$ esta a la \textbf{derecha} de $\vec{u}$ (giro horario / CW).
		\item $\vec{u} \times \vec{v} = 0 \iff \vec{u}$ y $\vec{v}$ son \textbf{colineales / paralelos}.
	\end{itemize}
\end{itemize}

\begin{center}
\begin{tikzpicture}[scale=0.8, >=Latex]
	% Ejes o base
	\draw[->, thick, blue!80!black] (0,0) -- (3,0) node[right] {$\vec{u}$};
	\draw[->, thick, red!80!black] (0,0) -- (2,2) node[above right] {$\vec{v}$ ($\text{CCW}: \vec{u} \times \vec{v} > 0$)};
	\draw[->, thick, green!50!black] (0,0) -- (2,-1.5) node[below right] {$\vec{w}$ ($\text{CW}: \vec{u} \times \vec{w} < 0$)};
	\draw[dashed, gray] (2,2) -- (2,0);
	\draw[thick, purple] (0,0) -- (2,0) node[midway, below] {\footnotesize $\text{proj}_{\vec{u}}(\vec{v})$};
	\draw (0.6,0) arc (0:45:0.6) node[midway, right=1pt] {\footnotesize $\theta$};
\end{tikzpicture}
\end{center}

\paragraph{2. Criterios de Paralelismo y Perpendicularidad.}
Dados dos vectores $\vec{u} = (x_1, y_1)$ y $\vec{v} = (x_2, y_2)$:
\begin{itemize}\setlength\itemsep{0pt}
	\item \textbf{Paralelos} ($\vec{u} \parallel \vec{v}$):
	\[
		\vec{u} \times \vec{v} = x_1 y_2 - y_1 x_2 = 0
	\]
	Misma direccion y sentido si $\vec{u} \cdot \vec{v} > 0$; sentidos opuestos si $\vec{u} \cdot \vec{v} < 0$.
	\item \textbf{Perpendiculares} ($\vec{u} \perp \vec{v}$):
	\[
		\vec{u} \cdot \vec{v} = x_1 x_2 + y_1 y_2 = 0
	\]
	\item \textbf{Vector perpendicular rotado}:
	Rotar $\vec{u} = (x, y)$ en $90^\circ$ antihorario da $\vec{u}^\perp = (-y, x)$. Rotar $-90^\circ$ (horario) da $(y, -x)$.
\end{itemize}

\paragraph{3. Orientacion de 3 puntos (Prueba CCW).}
Dados los puntos $A, B, C$, evaluar:
\[
	\text{orient}(A, B, C) = (B - A) \times (C - A) = (x_B - x_A)(y_C - y_A) - (y_B - y_A)(x_C - x_A)
\]
\begin{itemize}\setlength\itemsep{0pt}
	\item $> 0$: $C$ queda a la izquierda de la recta dirigida $A \to B$ (vuelta a la izquierda / CCW).
	\item $< 0$: $C$ queda a la derecha de la recta dirigida $A \to B$ (vuelta a la derecha / CW).
	\item $= 0$: $A, B, C$ son colineales.
\end{itemize}

\paragraph{4. Lineas y Distancias.}
\begin{itemize}\setlength\itemsep{0pt}
	\item \textbf{Ecuacion general de la recta}: $Ax + By + C = 0$.
	Dados dos puntos $P$ y $Q$: $A = y_Q - y_P$, $B = x_P - x_Q$, $C = -(A x_P + B y_P)$.
	\item \textbf{Distancia punto a recta}:
	Distancia del punto $P$ a la recta que pasa por $A$ y $B$:
	\[
		d(P, \text{recta}(AB)) = \frac{|(B - A) \times (P - A)|}{|B - A|} = \frac{|A x_P + B y_P + C|}{\sqrt{A^2 + B^2}}
	\]
	\item \textbf{Distancia punto a segmento $AB$}:
	\begin{itemize}\setlength\itemsep{0pt}
		\item Si $(B - A) \cdot (P - A) \le 0$, el punto mas cercano del segmento es $A \implies d = |P - A|$.
		\item Si $(A - B) \cdot (P - B) \le 0$, el punto mas cercano del segmento es $B \implies d = |P - B|$.
		\item En otro caso, la proyeccion cae dentro del segmento $\implies d = d(P, \text{recta}(AB))$.
	\end{itemize}
\end{itemize}

\begin{center}
\begin{tikzpicture}[scale=0.75, >=Latex]
	% Segmento AB
	\coordinate (A) at (0,0);
	\coordinate (B) at (4,0);
	\draw[ultra thick, black] (A) -- (B) node[midway, below] {Segmento $AB$};
	\fill (A) circle (2pt) node[below] {$A$};
	\fill (B) circle (2pt) node[below] {$B$};

	% Zonas
	\draw[dashed, gray] (0,-1) -- (0,2);
	\draw[dashed, gray] (4,-1) -- (4,2);

	% Puntos P1, P2, P3
	\coordinate (P1) at (-1,1.5);
	\coordinate (P2) at (2,1.8);
	\coordinate (P3) at (5,1.3);

	\fill[red] (P1) circle (2pt) node[above] {$P_1$};
	\draw[->, red, thick] (P1) -- (A) node[midway, left] {\footnotesize $|P_1 - A|$};

	\fill[blue] (P2) circle (2pt) node[above] {$P_2$};
	\draw[->, blue, thick] (P2) -- (2,0) node[midway, right] {\footnotesize $d(P_2, \text{recta})$};

	\fill[green!60!black] (P3) circle (2pt) node[above] {$P_3$};
	\draw[->, green!60!black, thick] (P3) -- (B) node[midway, right] {\footnotesize $|P_3 - B|$};
\end{tikzpicture}
\end{center}

\begin{itemize}\setlength\itemsep{0pt}
	\item \textbf{Interseccion de dos rectas}:
	$A_1 x + B_1 y + C_1 = 0$ y $A_2 x + B_2 y + C_2 = 0$.
	Determinante $D = A_1 B_2 - A_2 B_1$.
	\begin{itemize}\setlength\itemsep{0pt}
		\item Si $D = 0$: las rectas son paralelas ($C_1 A_2 \ne C_2 A_1$) o coincidentes ($C_1 A_2 = C_2 A_1$).
		\item Si $D \ne 0$: interseccion unica en:
		\[
			x = \frac{B_1 C_2 - B_2 C_1}{D}, \quad y = \frac{C_1 A_2 - C_2 A_1}{D}
		\]
	\end{itemize}
\end{itemize}

\paragraph{5. Circunferencias (Circulos).}
\begin{itemize}\setlength\itemsep{0pt}
	\item \textbf{Ecuacion}: $(x - h)^2 + (y - k)^2 = r^2$, con centro $(h, k)$ y radio $r$.
	\item \textbf{Interseccion recta-circunferencia}:
	Calcular $d = d(\text{centro}, \text{recta})$.
	\begin{itemize}\setlength\itemsep{0pt}
		\item $d > r$: 0 intersecciones (exterior).
		\item $d = r$: 1 interseccion (tangente en la proyeccion $P_{\text{proj}}$ del centro sobre la recta).
		\item $d < r$: 2 intersecciones. Los puntos estan a distancia $h = \sqrt{r^2 - d^2}$ a ambos lados de $P_{\text{proj}}$ a lo largo del vector unitario director de la recta.
	\end{itemize}
	\item \textbf{Interseccion entre dos circunferencias}:
	Centros $C_1, C_2$ con radios $r_1, r_2$ y distancia $d = |C_1 - C_2|$:
	\begin{itemize}\setlength\itemsep{0pt}
		\item $d > r_1 + r_2$ o $d < |r_1 - r_2|$: no se intersectan (disjuntos externos o uno contenido en otro).
		\item $d = 0$ y $r_1 = r_2$: infinitos puntos (circunferencias identicas).
		\item $d = r_1 + r_2$ o $d = |r_1 - r_2|$: 1 punto (tangencia exterior o interior).
		\item En otro caso: 2 puntos de corte simetricos respecto a la linea de centros.
	\end{itemize}
	\item \textbf{Circuncentro de un triangulo $\triangle ABC$}:
	Punto de interseccion de las tres mediatrices.
	Radio de la circunferencia circunscrita:
	\[
		R = \frac{a \cdot b \cdot c}{4 \cdot \text{Area}} = \frac{|B - A| \, |C - B| \, |A - C|}{2 \, |(B - A) \times (C - A)|}
	\]
\end{itemize}

\begin{center}
\begin{tikzpicture}[scale=0.75, >=Latex]
	% Circunferencia y recta secante
	\coordinate (O) at (0,0);
	\draw[thick, blue] (O) circle (1.5);
	\fill (O) circle (2pt) node[below left] {Centro};

	\draw[thick, red!80!black] (-2, 0.9) -- (2.5, 0.9) node[right] {Recta};
	\draw[dashed, gray] (O) -- (0, 0.9) node[midway, left] {$d$};
	\draw[dashed, blue] (O) -- (1.2, 0.9) node[midway, below right] {$r$};
	\fill (0, 0.9) circle (1.5pt) node[above] {$P_{\text{proj}}$};
	\fill[black] (-1.2, 0.9) circle (2pt);
	\fill[black] (1.2, 0.9) circle (2pt);
	\draw[<->, purple] (-1.2, 1.15) -- (0, 1.15) node[midway, above] {\footnotesize $\sqrt{r^2 - d^2}$};
\end{tikzpicture}
\end{center}

\paragraph{6. Tips criticos de precision y tipos en CP.}
\begin{itemize}\setlength\itemsep{0pt}
	\item \textbf{Evitar flotantes en orientacion}: Al hacer cross product de coordenadas $\le 10^9$, el resultado puede llegar a $2 \cdot 10^{18}$, cabiendo en \texttt{long long}. Si se multiplica con otras distancias, usar \verb|__int128_t|.
	\item \textbf{Comparacion con EPS}: Para flotantes (\texttt{double}), usar constante $\varepsilon = 10^{-9}$:
	\begin{itemize}\setlength\itemsep{0pt}
		\item Igualdad: \verb|abs(a - b) < EPS|
		\item Menor estricto: \verb|a < b - EPS|
		\item Menor o igual: \verb|a < b + EPS|
	\end{itemize}
	\item \textbf{Angulo de un vector}: Usar \verb|atan2(y, x)|, que devuelve el angulo exacto en $[-\pi, \pi]$ considerando signos de ambos cuadrantes sin division por cero.
\end{itemize}
\vspace{0.5em}
